% Register bank conflicts in a 4-bank file with 2 read ports per bank
% (Lesson 4).  (a) reading R0 of T0..T3 touches one entry in every bank:
% one cycle.  (b) reading R0..R3 of T0 touches four entries of bank 0:
% two cycles.  Redraws the outline figures m4l2_bank_conflict2 and
% m4l2_bank_conflict3.
% Usage: % Register bank conflicts in a 4-bank file with 2 read ports per bank
% (Lesson 4).  (a) reading R0 of T0..T3 touches one entry in every bank:
% one cycle.  (b) reading R0..R3 of T0 touches four entries of bank 0:
% two cycles.  Redraws the outline figures m4l2_bank_conflict2 and
% m4l2_bank_conflict3.
% Usage: % Register bank conflicts in a 4-bank file with 2 read ports per bank
% (Lesson 4).  (a) reading R0 of T0..T3 touches one entry in every bank:
% one cycle.  (b) reading R0..R3 of T0 touches four entries of bank 0:
% two cycles.  Redraws the outline figures m4l2_bank_conflict2 and
% m4l2_bank_conflict3.
% Usage: % Register bank conflicts in a 4-bank file with 2 read ports per bank
% (Lesson 4).  (a) reading R0 of T0..T3 touches one entry in every bank:
% one cycle.  (b) reading R0..R3 of T0 touches four entries of bank 0:
% two cycles.  Redraws the outline figures m4l2_bank_conflict2 and
% m4l2_bank_conflict3.
% Usage: \input{figures/l04-register-bank-conflict}
% Width: body (120 mm); the drawing is about 118 mm wide.
\begin{tikzpicture}[x=1mm, y=1mm, font=\scriptsize\sffamily]
  % \bankfile{x}{title}{highlight style}: 4 banks x 4 entries
  \newcommand{\bankfile}[2]{%
    \node[anchor=south west, font=\scriptsize\sffamily\bfseries] at (#1,24) {#2};
    \foreach \b in {0,...,3} {
      \node[anchor=south, text=ln-muted] at (#1+\b*12+5,18.5) {\iflangja{バンク \b}{bank \b}};
      \foreach \r in {0,...,3} {
        \draw[fill=ln-box] (#1+\b*12,18-\r*4.2) rectangle ++(10,-4) node[midway] {T\b.R\r};
      }
    }}
  \bankfile{0}{\iflangja{(a) T0--T3 の R0：1 サイクル}{(a) R0 of T0--T3: one cycle}}
  \foreach \b in {0,...,3} {
    \draw[ln-tip, very thick] (\b*12,18) rectangle ++(10,-4);
  }
  \node[anchor=north west, align=left, text=ln-muted] at (0,-0.5)
    {\iflangja{バンクごとに 1 読み出し}{one read per bank}};
  \bankfile{62}{\iflangja{(b) T0 の R0--R3：2 サイクル}{(b) R0--R3 of T0: two cycles}}
  \foreach \r in {0,...,3} {
    \draw[ln-caution, very thick] (62,18-\r*4.2) rectangle ++(10,-4);
  }
  \node[anchor=east, text=ln-caution, align=right] at (61,14.2) {\iflangja{1 周目}{cycle 1}};
  \node[anchor=east, text=ln-caution, align=right] at (61,5.8) {\iflangja{2 周目}{cycle 2}};
  \draw[ln-caution] (61.2,12.1) -- (61.2,18.2);
  \draw[ln-caution] (61.2,3.7) -- (61.2,9.8);
  \node[anchor=north west, align=left, text=ln-muted] at (62,-0.5)
    {\iflangja{バンク 0 に 4 読み出し，ポートは 2 本}{four reads to bank 0, two ports}};
\end{tikzpicture}

% Width: body (120 mm); the drawing is about 118 mm wide.
\begin{tikzpicture}[x=1mm, y=1mm, font=\scriptsize\sffamily]
  % \bankfile{x}{title}{highlight style}: 4 banks x 4 entries
  \newcommand{\bankfile}[2]{%
    \node[anchor=south west, font=\scriptsize\sffamily\bfseries] at (#1,24) {#2};
    \foreach \b in {0,...,3} {
      \node[anchor=south, text=ln-muted] at (#1+\b*12+5,18.5) {\iflangja{バンク \b}{bank \b}};
      \foreach \r in {0,...,3} {
        \draw[fill=ln-box] (#1+\b*12,18-\r*4.2) rectangle ++(10,-4) node[midway] {T\b.R\r};
      }
    }}
  \bankfile{0}{\iflangja{(a) T0--T3 の R0：1 サイクル}{(a) R0 of T0--T3: one cycle}}
  \foreach \b in {0,...,3} {
    \draw[ln-tip, very thick] (\b*12,18) rectangle ++(10,-4);
  }
  \node[anchor=north west, align=left, text=ln-muted] at (0,-0.5)
    {\iflangja{バンクごとに 1 読み出し}{one read per bank}};
  \bankfile{62}{\iflangja{(b) T0 の R0--R3：2 サイクル}{(b) R0--R3 of T0: two cycles}}
  \foreach \r in {0,...,3} {
    \draw[ln-caution, very thick] (62,18-\r*4.2) rectangle ++(10,-4);
  }
  \node[anchor=east, text=ln-caution, align=right] at (61,14.2) {\iflangja{1 周目}{cycle 1}};
  \node[anchor=east, text=ln-caution, align=right] at (61,5.8) {\iflangja{2 周目}{cycle 2}};
  \draw[ln-caution] (61.2,12.1) -- (61.2,18.2);
  \draw[ln-caution] (61.2,3.7) -- (61.2,9.8);
  \node[anchor=north west, align=left, text=ln-muted] at (62,-0.5)
    {\iflangja{バンク 0 に 4 読み出し，ポートは 2 本}{four reads to bank 0, two ports}};
\end{tikzpicture}

% Width: body (120 mm); the drawing is about 118 mm wide.
\begin{tikzpicture}[x=1mm, y=1mm, font=\scriptsize\sffamily]
  % \bankfile{x}{title}{highlight style}: 4 banks x 4 entries
  \newcommand{\bankfile}[2]{%
    \node[anchor=south west, font=\scriptsize\sffamily\bfseries] at (#1,24) {#2};
    \foreach \b in {0,...,3} {
      \node[anchor=south, text=ln-muted] at (#1+\b*12+5,18.5) {\iflangja{バンク \b}{bank \b}};
      \foreach \r in {0,...,3} {
        \draw[fill=ln-box] (#1+\b*12,18-\r*4.2) rectangle ++(10,-4) node[midway] {T\b.R\r};
      }
    }}
  \bankfile{0}{\iflangja{(a) T0--T3 の R0：1 サイクル}{(a) R0 of T0--T3: one cycle}}
  \foreach \b in {0,...,3} {
    \draw[ln-tip, very thick] (\b*12,18) rectangle ++(10,-4);
  }
  \node[anchor=north west, align=left, text=ln-muted] at (0,-0.5)
    {\iflangja{バンクごとに 1 読み出し}{one read per bank}};
  \bankfile{62}{\iflangja{(b) T0 の R0--R3：2 サイクル}{(b) R0--R3 of T0: two cycles}}
  \foreach \r in {0,...,3} {
    \draw[ln-caution, very thick] (62,18-\r*4.2) rectangle ++(10,-4);
  }
  \node[anchor=east, text=ln-caution, align=right] at (61,14.2) {\iflangja{1 周目}{cycle 1}};
  \node[anchor=east, text=ln-caution, align=right] at (61,5.8) {\iflangja{2 周目}{cycle 2}};
  \draw[ln-caution] (61.2,12.1) -- (61.2,18.2);
  \draw[ln-caution] (61.2,3.7) -- (61.2,9.8);
  \node[anchor=north west, align=left, text=ln-muted] at (62,-0.5)
    {\iflangja{バンク 0 に 4 読み出し，ポートは 2 本}{four reads to bank 0, two ports}};
\end{tikzpicture}

% Width: body (120 mm); the drawing is about 118 mm wide.
\begin{tikzpicture}[x=1mm, y=1mm, font=\scriptsize\sffamily]
  % \bankfile{x}{title}{highlight style}: 4 banks x 4 entries
  \newcommand{\bankfile}[2]{%
    \node[anchor=south west, font=\scriptsize\sffamily\bfseries] at (#1,24) {#2};
    \foreach \b in {0,...,3} {
      \node[anchor=south, text=ln-muted] at (#1+\b*12+5,18.5) {\iflangja{バンク \b}{bank \b}};
      \foreach \r in {0,...,3} {
        \draw[fill=ln-box] (#1+\b*12,18-\r*4.2) rectangle ++(10,-4) node[midway] {T\b.R\r};
      }
    }}
  \bankfile{0}{\iflangja{(a) T0--T3 の R0：1 サイクル}{(a) R0 of T0--T3: one cycle}}
  \foreach \b in {0,...,3} {
    \draw[ln-tip, very thick] (\b*12,18) rectangle ++(10,-4);
  }
  \node[anchor=north west, align=left, text=ln-muted] at (0,-0.5)
    {\iflangja{バンクごとに 1 読み出し}{one read per bank}};
  \bankfile{62}{\iflangja{(b) T0 の R0--R3：2 サイクル}{(b) R0--R3 of T0: two cycles}}
  \foreach \r in {0,...,3} {
    \draw[ln-caution, very thick] (62,18-\r*4.2) rectangle ++(10,-4);
  }
  \node[anchor=east, text=ln-caution, align=right] at (61,14.2) {\iflangja{1 周目}{cycle 1}};
  \node[anchor=east, text=ln-caution, align=right] at (61,5.8) {\iflangja{2 周目}{cycle 2}};
  \draw[ln-caution] (61.2,12.1) -- (61.2,18.2);
  \draw[ln-caution] (61.2,3.7) -- (61.2,9.8);
  \node[anchor=north west, align=left, text=ln-muted] at (62,-0.5)
    {\iflangja{バンク 0 に 4 読み出し，ポートは 2 本}{four reads to bank 0, two ports}};
\end{tikzpicture}
